% ============================================================
%  ch02  第一讲 §1.3：附着于去心圆盘的群
%  原文行 909–1032
% ============================================================

\subsection{附着于去心圆盘的群}

$\Gr_{\underline G}$ 还有一种有用的解释。先引入射流群与环路群。

\begin{dfn}
设 $X$ 是 $\mO=k[[t]]$ 上的预层。$X$ 的 \textbf{$n$-射流空间} $L^nX$ 是这样一个预层：它给每个 $R$ 配上集合
\[
L^nX(R)=X(R[t]/t^n).
\]
$X$ 的\textbf{形式射流空间}（有时也称为\textbf{弧空间}或\textbf{正环路空间}）$L^+X$ 是这样一个预层：它给每个 $R$ 配上集合
\[
L^+X(R)=X(R[[t]]).
\]
设 $X$ 是 $F=k((t))$ 上的预层。$X$ 的\textbf{环路空间} $LX$ 是这样一个预层：它给每个 $R$ 配上集合
\[
LX(R)=X(R((t))).
\]
若 $X$ 是 $k$ 上的预层，为简单起见，在不引起混淆时把 $L^+(X\otimes_k\mO)$ 记作 $L^+X$，把 $L(X\otimes_k F)$ 记作 $LX$。
\end{dfn}

\begin{prop}~
\begin{enumerate}
\item 设 $X$ 是 $\mO=k[[t]]$ 上有限型概形。则 $L^nX$ 可由 $k$ 上有限型概形表示。此外
\[
L^+X\cong \varprojlim\ L^nX
\]
可由概形表示。若 $X$ 仿射，则 $L^+X$ 也仿射。若 $f:X\to Y$ 是开（相应地，闭）嵌入，则 $L^+f: L^+X\to L^+Y$ 也是。

\item 设 $X$ 是 $F=k((t))$ 上有限型仿射概形。则 $LX$ 可由 $k$ 上归纳仿射概形表示。若 $f:X\to Y$ 是闭嵌入，则 $Lf:LX\to LY$ 也是。此外，若 $X=\underline X\otimes_\mO F$，其中 $\underline X$ 是 $\mO$ 上有限型仿射概形，则 $L^+\underline X\subset LX$ 是闭子概形。
\end{enumerate}
\end{prop}

\begin{proof}见 \cite[1.a]{PR}。
\end{proof}

\begin{ex}
若 $X=\bA^1$，则 $L^+X(R)=R[[t]]=\{\sum r_it^i \mid r_i\in R,\ i<0 \text{ 时 } r_i=0\}$，而 $LX(R)=R((t))=\{\sum r_it^i\mid  r_i\in R,\ i\ll 0 \text{ 时 } r_i=0\}$。因此 $L^+X\cong \Spec k[r_0,r_1,\ldots]$，而 $LX=\varinjlim_i \Spec k[r_{-i}, r_{-i+1},\ldots]$。
\end{ex}

\begin{ex}\label{hyperspacial and Iwahori2}
考虑例 \ref{hyperspacial and Iwahori} 中的群概形。自然同态 $\underline G'\to \underline G$ 诱导闭嵌入 $L^+\underline G'\subset L^+\underline G$。事实上，求值映射可以提升为仿射 $k$-群的映射 $\on{ev}:L^+G\to G$，而 $L^+\underline G'=\on{ev}^{-1}(B)$。但要注意，同态 $\underline G'\to \underline G$ 并不是闭嵌入，因为基变换到 $F$ 之后它是同构。
\end{ex}

回到仿射 Grassmann 簇。注意有自然作用
\begin{equation}\label{I:action}
L\underline G\times \Gr_{\underline G}\to \Gr_{\underline G},\quad (A, (\mE,\beta))\mapsto (\mE,A\beta).
\end{equation}

\begin{prop}\label{I:quotient}
仿射 Grassmann 簇 $\Gr_{\underline G}$ 可以等同于 fpqc 商 $[L\underline G/L^+\underline G]$。
\end{prop}

\begin{proof}
先需要一个引理。

\begin{lem}\label{et loc triv}
$D_R$ 上每个 $\underline G$-主丛都能在某个 \'etale 覆盖 $\Spec R'\to \Spec R$ 的 $D_{R'}$ 上平凡化。
\end{lem}

\begin{proof}
首先，映射 $\mE\otimes_{R[[t]]}R\to\Spec R$ 光滑（因为 $\underline G$ 光滑），故在某个 \'etale 覆盖 $\Spec R'\to \Spec R$ 上有截面。又由 $\mE\to D_R$ 的光滑性，该截面可提升为 $\mE$ 在 $D_{R'}$ 上的截面 $\epsilon$。
\end{proof}

现在来证命题。要证的论断等价于说 $L\underline G$ 是 $\Gr_{\underline G}$ 上的 $L^+\underline G$-主丛。确切地说，对 $k$-代数 $R$ 有典范同构
\[
L\underline G(R)\cong\left\{(\mE,\beta, \epsilon)\ \left|\ \begin{array}{l}
(\mE,\beta)\in \Gr_{\underline G}(R),\\
\epsilon: \mE^0\cong \mE \text{ 是平凡化}
\end{array}\right.\right\}.
\]
事实上，所要的同构由 $A\mapsto (\mE^0,A,\on{id})$ 给出，其逆映射为 $(\mE,\beta,\epsilon)\mapsto A:=\beta\epsilon$。
\end{proof}

\begin{rmk}
由命题的证明，$L^+\underline G$-主丛 $L\underline G\to \Gr_{\underline G}$ 在 \'etale 拓扑下局部平凡。若 $\underline{G}=G\otimes \mO$ 是常值约化的（更一般地，是所谓驯顺分歧的），它在 Zariski 拓扑下也局部平凡（见引理 \ref{big cell}）。
\end{rmk}

卷积 Grassmann 簇 \eqref{I:conv} 也有类似的解释。利用 $L^+\underline G$-主丛 $L\underline G\to \Gr_{\underline G}$ 以及 $L^+\underline G$ 在 $\Gr_{\underline G}$ 上的左作用（由 \eqref{I:action} 定义），可作扭曲积 $L\underline G\times^{L^+ \underline G}\Gr_{\underline G}$（见 \S\ref{Intro:torsor}）。同样的论证表明 $L\underline G\times^{L^+ \underline G}\Gr_{\underline G}$ 同构于 $\Gr_{\underline G}\tilde\times\Gr_{\underline G}$（特别地，我们的记号是相容的）。

由此得到一个直接推论。

\begin{cor}
仿射 Grassmann 簇 $\Gr_{\underline G}$ 形式光滑。
\end{cor}

\begin{proof}
事实上，$L\underline G$ 与 $L^+\underline G$ 都形式光滑。
\end{proof}

\begin{rmk}\label{formal smoothness}
对归纳概形而言，形式光滑是个很弱的概念。例如 $L\bG_m$ 形式光滑，但容易看出它\textbf{极不}既约。
\end{rmk}

下面是环路群的一个更深入的应用。设 $k$ 代数闭，$\underline G_F=\underline G\otimes F$ 连通约化。记 $I=\Gal(\overline F/F)$，它作用在 $\underline G_F$ 的基本群 $\pi_1(\underline G_F)$ 上。以 $\pi_1(\underline G_F)_I$ 记其余不变量群。以 $p$ 记 $k$ 的特征指数。

\begin{thm}\label{components}~
\begin{enumerate}
\item 有典范同构 $\pi_0(L\underline G)\cong \pi_1(\underline G_F)_I$。映射 $L\underline G\to \Gr_{\underline G}$ 诱导同构 $\pi_0(L\underline G)\cong \pi_0(\Gr_{\underline G})$。以 $(\Gr_{\underline G})^0$ 记单位连通分支。
\item 设 $\underline G_F$ 半单且单连通。则 $\Gr_{\underline G}$ 既约。
\item 设 $\underline G_F$ 半单且 $p\nmid \pi_1(\underline G_F)$。以 $G_{\on{sc}}$ 记 $\underline G_F$ 的单连通覆盖，并设它可延拓为 \'etale 覆盖 $\underline G_{\on{sc}}\to \underline G$。则它诱导同构 $\Gr_{\underline G_{\on{sc}}}\cong (\Gr_{\underline G})^0$。特别地，$\Gr_{\underline G}$ 既约。
\item 一般地，设 $p\nmid \pi_1(G_\der)$，其中 $G_\der$ 是 $\underline G_F$ 的导群。则映射 $\underline G_{\on{sc}}\to \underline G$ 诱导同构 $\Gr_{\underline G_{\on{sc}}}\cong (\Gr_{\underline G})^0_{\on{red}}$。
\end{enumerate}
\end{thm}

这个定理（在各种一般性下）的证明可从多处提取（\cite[\S 4.5]{BD}、\cite{BL}、\cite{LS}、\cite{BLS}、\cite{Fa}，最终可追溯到 \cite{PR}）。注意 (1) 类似于拓扑中众所周知的事实：环路空间的 $\pi_0$ 等于原空间的 $\pi_1$（当 $k=\bC$ 时这其实不止是类比，见定理 \ref{top triv}）。

由该定理，$\Gr_{\underline G}$ 上的许多问题都可化归到 $G$ 单连通的情形。

\begin{ex}
(i) 取 $G=\bG_m$。由定理，$\pi_0(\Gr_G)\cong \bZ$。事实上 $(\Gr_G)_{\on{red}}=\pi_0(G)=\bZ$。但如注 \ref{formal smoothness} 所述，它不既约。

(ii) 在 \cite[Remark 6.4]{PR} 中证明了：当 $G=\on{PGL}_2$、$\cha\ k=2$ 时，$\Gr_G$ 不既约。
\end{ex}

除 $L\underline G$ 在 $\Gr_{\underline G}$ 上的自然作用外，$\Gr_{\underline G}$ 上还有另一个重要的对称性——至少在 $\underline G=G\otimes_k\mO$ 为常值时如此。以 $\Aut(D)$ 记 $\mO=k[[t]]$ 的自同构群归纳概形。确切地说，它是 $\Aff_k$ 上的群预层，定义为
\[
\Aut(D)(R)=\Aut_R(R[[t]]).
\]
$R[[t]]$ 的任一 $R$-自同构 $\varphi$ 由 $\varphi(t)=a_0+a_1t+a_2t^2+\cdots\in R[[t]]$ 的像决定。容易验证：存在 $\psi$ 使 $\psi(\varphi(t))=t$，当且仅当 $a_0$ 幂零且 $a_1$ 可逆。因此
\[
\Aut(D)\cong \on{Spf} k[[a_0]] \times \Spec k[a_1^{\pm 1}]\times \Spec\ k[a_2,\cdots].
\]
注意 $\Aut^+(D)=\Spec k[a_1^{\pm 1},a_2,\cdots]$ 是闭正规子群概形，它分类那些模 $t$ 后为恒等映射的 $R[[t]]$ 自同构。于是
\begin{equation}\label{decomp1}
\Aut(D)\cong \Aut^+(D)\rtimes \hat{\bG}_a,
\end{equation}
其中 $\hat{\bG}_a=\on{Spf} k[[a_0]]$ 是形式加法群。此外，$\Aut^{++}(D)=\Spec\ k[a_2,\cdots]$ 是 $\Aut^+(D)$ 的超幂幺根。记 $\bG_m^{\on{rot}}=\Spec k[a_1^{\pm 1}]$，通常称为\textbf{旋转环面}。则
\begin{equation}\label{decomp2}
\Aut^+(D)=\Aut^{++}(D)\rtimes \bG_m^{\on{rot}}.
\end{equation}

\begin{rmk}
若 $F$ 是一般的局部域，且没有指定的素元可选，则不存在典范的半直积分解 \eqref{decomp1} 与 \eqref{decomp2}。但相应的群归纳概形短正合列仍然有定义。
\end{rmk}

现在，若 $\underline G$ 是 $k[[t]]$ 上的群概形，带有提升 $k[[t]]$ 上典范作用的 $\Aut(D)$-作用（例如常值的 $\underline G=G\otimes_k k[[t]]$），则有作用
\begin{equation}\label{auto of disc}
\Aut(D)\times \Gr_{\underline G}\to \Gr_{\underline G},\quad (g, (\mE,\beta))\mapsto (g^*\mE,g^*\beta).
\end{equation}
特别地，旋转环面 $\bG_m^{\on{rot}}$ 作用在 $\Gr_{\underline G}$ 上。
